% =============================================================================
% 07_deployment.tex — lane WC
% Sources: .github/workflows/ci.yml, ci-hardware.yml, flake.nix, flake.lock,
%          devenv.yaml, docs/remote-box.md, docs/upstream_contribution.md,
%          upstream/tt-mlir/README.md, scripts/build-tt-mlir.sh
% =============================================================================
\section{Deployment}\label{sec:deployment}

Deployment, for TensixFlow, means reproducibility of the toolchain plus two
CI paths: an always-on CPU path and a double-gated hardware path. There is
no cloud provisioning anywhere in the loop.

\subsection{Reproducibility: the pinned nix toolchain}

The development environment is a nix flake
(\texttt{flake.nix} + committed \texttt{flake.lock}):
\begin{itemize}[nosep]
  \item \textbf{Pinned input.} \texttt{flake.lock} locks
        \texttt{NixOS/nixpkgs} to rev
        \texttt{bd3c09ca2a2858daa15f26eb}\allowbreak\texttt{557faa4709635a41}
        (\texttt{nixos-unstable}), so every shell resolves the same package
        set on every machine.
  \item \textbf{Pinned LLVM.} \texttt{flake.nix} selects the explicit
        \texttt{llvmPackages\_20} set --- LLVM/MLIR 20.1.8 --- and its
        comment forbids the floating default \texttt{llvmPackages}. The
        shell exports \texttt{LLVM\_DIR}/\texttt{MLIR\_DIR} from the pinned
        store paths and puts \texttt{mlir-tblgen} on the PATH; the package
        list also fixes cmake, ninja, lit, and GoogleTest (with the
        \texttt{dev} output named explicitly so \texttt{find\_package(GTest)}
        resolves).
  \item \textbf{devenv layering.} \texttt{devenv.yaml} re-pins its
        \texttt{nixpkgs} input to the \emph{same} rev
        (\allowbreak\texttt{\scriptsize bd3c09ca2a2858daa15f26eb}\allowbreak\texttt{\scriptsize 557faa4709635a41}) with an explicit
        warning: at a floating \texttt{nixos-unstable} URL,
        \texttt{devenv shell} re-resolves the lock to a head whose mlir is
        uncached and source-broken. The rev must stay pinned here, not only
        in \texttt{devenv.lock}.
  \item \textbf{The lit pin.} Configure passes
        \texttt{-DLLVM\_EXTERNAL\_LIT=\$(dirname \$(which lit))/.lit-wrapped}:
        under nix, \texttt{bin/lit} is a bash wrapper, and LLVM's
        \texttt{add\_lit\_target} spawns lit as \texttt{python3 <lit>}, which
        breaks on the wrapper. Pointing CMake at the real Python script
        \texttt{.lit-wrapped} keeps the lit target working.
\end{itemize}

The exact configure command, as documented in the README:

\begin{lstlisting}[language=bash]
cmake -G Ninja -B build \
  -DMLIR_DIR=$MLIR_DIR -DLLVM_DIR=$LLVM_DIR \
  -DLLVM_EXTERNAL_LIT="$(dirname "$(which lit)")/.lit-wrapped"
\end{lstlisting}

One sentence on this document itself: the \LaTeX{} report builds in an
ephemeral \texttt{nix shell nixpkgs\#texlive.combined.scheme-full} (TeX Live
2025) that resolves through the same pinned registry input, without touching
the dev flake --- see \texttt{report/BUILDING.md}.

\subsection{CI topology}\label{sec:deployment-ci}

Two workflows, two disjoint runner classes
(Figure~\ref{fig:deploy-ci-topology}).

\begin{figure}[htbp]
  \centering
  \begin{tikzpicture}[
      box/.style={draw=tensixnavy, rounded corners=2pt, align=center,
                  font=\scriptsize, inner sep=3pt, fill=white,
                  minimum width=6.4cm},
      gate/.style={box, fill=tensixamber!25},
      ok/.style={box, fill=tensixteal!15},
      head/.style={box, fill=tensixlight, font=\small},
      arr/.style={-{Stealth[length=2.2mm]}, tensixnavy},
      paneltitle/.style={font=\small\itshape, tensixnavy},
    ]
    % ---- left: CPU path --------------------------------------------------------
    \node[head] (c0) at (-3.9,0) {\textbf{\texttt{ci.yml} --- CPU path}\\
      \scriptsize on: push, pull\_request (every push)};
    \node[box, below=0.36cm of c0] (c1) {runner: \texttt{ubuntu-latest}\\
      nix via \texttt{cachix/install-nix-action@v31} (flakes enabled)};
    \node[box, below=0.3cm of c1] (c2) {\texttt{nix develop} configure + build\\
      \scriptsize incl.\ the \texttt{.lit-wrapped} LLVM\_EXTERNAL\_LIT pin};
    \node[box, below=0.3cm of c2] (c3) {lit regression: \texttt{check-tensixflow}\\
      \scriptsize 25 tests = 24 passed + 1 unsupported};
    \node[box, below=0.3cm of c3] (c4) {ctest: 11 gtest units\\
      \scriptsize incl.\ \texttt{hw-acceptance} $\Rightarrow$ generated-kernel\\
      \scriptsize compile check (no-device branch)};
    \node[box, below=0.3cm of c4] (c5) {pytest: golden / harness / DSL suites};
    \draw[arr] (c0) -- (c1); \draw[arr] (c1) -- (c2); \draw[arr] (c2) -- (c3);
    \draw[arr] (c3) -- (c4); \draw[arr] (c4) -- (c5);
    \begin{scope}[on background layer]
      \node[fit=(c0)(c5), draw=tensixteal, rounded corners=3pt,
            fill=tensixteal!6, inner sep=7pt] (cpupanel) {};
    \end{scope}
    \node[paneltitle, anchor=south] at (cpupanel.south)
      {always runs --- no device, no cloud spend};

    % ---- right: gated hardware path ---------------------------------------------
    \node[head] (h0) at (3.9,0) {\textbf{\texttt{ci-hardware.yml} --- gated path}\\
      \scriptsize on: pull\_request (label events), workflow\_dispatch};
    \node[gate, below=0.36cm of h0] (h1) {\textbf{Gate 1 (job \texttt{if}):}\\
      PR label \texttt{hardware.present}\\ OR explicit \texttt{workflow\_dispatch}};
    \node[gate, below=0.3cm of h1] (h2) {\textbf{Gate 2 (\texttt{runs-on}):}\\
      \texttt{[self-hosted, tenstorrent,}\\ \texttt{wormhole-n150]} only};
    \node[box, below=0.3cm of h2] (h3) {assert \texttt{/dev/tenstorrent} exists\\
      \scriptsize absent $\Rightarrow$ \texttt{exit 1}, never fall through\\
      \scriptsize (a deviceless runner is a misconfiguration)};
    \node[box, below=0.3cm of h3] (h4) {on-device branch only (no fallback):\\
      \scriptsize compile vs real tt-metal, run,\\
      \scriptsize PCC $\ge 0.9875$ vs the golden};
    \draw[arr] (h0) -- (h1); \draw[arr] (h1) -- (h2); \draw[arr] (h2) -- (h3);
    \draw[arr] (h3) -- (h4);
    \begin{scope}[on background layer]
      \node[fit=(h0)(h4), draw=tensixamber, rounded corners=3pt,
            fill=tensixamber!6, inner sep=7pt] (hwpanel) {};
    \end{scope}
    \node[paneltitle, anchor=south] at (hwpanel.south)
      {runs only where hardware exists --- fails loudly, never silently passes};
  \end{tikzpicture}
  \caption{CI topology: the CPU path (\texttt{ci.yml}) runs the full
  build/lit/ctest/pytest chain on every push, with the generated-kernel
  compile check riding inside ctest's \texttt{hw-acceptance} test; the
  hardware path (\texttt{ci-hardware.yml}) is double-gated
  (Section~\ref{sec:hardware-acceptance}) and only ever runs on a tagged
  self-hosted runner carrying the device.}
  \label{fig:deploy-ci-topology}
\end{figure}

Notes on the CPU path (all from \texttt{.github/workflows/ci.yml}): the
ctest step passes \texttt{--no-tests=ignore} and the pytest step tolerates
exit code 5 (``no tests ran'') so both stayed green in early waves; today
both carry real suites. The generated-kernel compile check is not a separate
CI step --- it is enforced inside ctest via \texttt{hw-acceptance}, whose
no-device branch runs
\path{scripts/hw_bringup_dryrun.sh && cmake --build build --target
check-generated-kernels-compile} (Section~\ref{sec:hardware-acceptance}).

\subsection{Cloud policy: opt-in, documented, never invoked}

The MVP is local-first by policy:
\begin{itemize}[nosep]
  \item \texttt{docs/remote-box.md} states it up front: ``You do not need
        this.'' A remote Linux build box is \emph{purely opt-in}, for keeping
        a laptop free; the document's provider notes (e.g.\ PrimeIntellect
        CPU sandboxes as the cheapest option) are indicative price ordering,
        to be verified in-console.
  \item \texttt{scripts/remote-box.sh} is an idempotent setup script run
        \emph{manually} on a machine the user already owns access to:
        install nix if absent, clone the repo, enter the dev shell, run the
        standard configure + build. It provisions nothing and is never
        invoked by CI.
  \item Nothing in \texttt{ci.yml}, \texttt{ci-hardware.yml}, or the test
        suites provisions or contacts a cloud provider, and nothing incurs
        cloud spend automatically. Likewise \texttt{scripts/build-tt-mlir.sh}
        (the optional tt-mlir study helper) is off the critical path: no
        CMake target depends on it.
\end{itemize}

\subsection{Upstream contribution: scouted, implemented locally, not submitted}

Status, exactly: TensixFlow \emph{scouted} a concrete tt-mlir contribution
and \emph{implemented it locally}; the submission process is fully
documented; and \textbf{nothing was submitted, forked, pushed, or merged}
--- per the project owner's explicit binding instruction. This is policy in
action, not an omission:

\begin{itemize}[nosep]
  \item \textbf{Scouted} (\texttt{docs/upstream\_contribution.md}, todo 58):
        one new lit test for tt-mlir's \texttt{ttir.softmax} verifier,
        covering both error branches --- one of which (shape mismatch) had
        \emph{zero} test coverage anywhere in tt-mlir at the scouted
        revision (\texttt{main} @ \texttt{da6f1ce0}). The gap is cited from
        the actual repo: verifier source, op definition, tree listings, and
        a search over the diagnostic strings.
  \item \textbf{Implemented locally} (\texttt{upstream/tt-mlir/}, todo 59):
        \texttt{softmax\_tests\_negative.mlir} --- 22 lines, three negative
        cases, CHECK strings quoted verbatim from \texttt{SoftmaxOp::verify()}
        --- plus a README recording provenance (scratch clone, \emph{not} a
        fork; branch \texttt{softmax-negative-tests} based on
        \texttt{da6f1ce0}; upstream lint hooks run; scoped diff of exactly
        one new file) and an honest ``what is NOT verified'' section (never
        executed against a real \texttt{ttmlir-opt}, whose toolchain is not
        built on this box).
  \item \textbf{Submission documented, never executed}:
        \texttt{docs/upstream\_contribution.md} records the complete process
        (fork + branch, \texttt{pre-commit}, targeted \texttt{llvm-lit},
        PR template, CI expectations) under a heading that says
        ``DOCUMENTED ONLY; never executed in this plan,'' and states that
        the plan ``explicitly stops before step 1.'' The same constraint is
        repeated in \texttt{upstream/tt-mlir/README.md}.
\end{itemize}
